% GENERATED FILE. Edit canonical lessons, then run npm run generate:books.
% canonical-source-hash: fnv1a64:83308bff44a48ca3
% canonical-lessons: IT-C503-scacchi, IT-C503-dama, IT-C503-dado, IT-C503-gioco-da-tavolo, IT-C503-barare

\chapter{Giochi da tavolo}
\label{ch:it-giochi-da-tavolo}
\hlchaptermodality{503}

\begin{chapteropening}
\textbf{By the end of this chapter:} \emph{I can talk about board games.}

Complete the last lesson of chapter 503: I can talk about board games.
\end{chapteropening}

\section[gli scacchi]{gli scacchi — chess}

\label{lesson:IT-C503-scacchi}



\begin{quote}
Before the new one: say the Italian for trip, then the Italian for to explore.
\end{quote}

\subsection*{You'll want to know: gli scacchi}
\textbf{gli scacchi} — \textquotedblleft{}chess\textquotedblright{}.

\begin{quote}
Mio nonno mi ha insegnato a giocare a scacchi.
\end{quote}

\begin{quote}
\emph{My grandfather taught me to play chess.}
\end{quote}

\emph{Giocare a scacchi}; \emph{scacco matto} — checkmate.

\begin{cousinweb}
\emph{scacchi} is the plural of \emph{scacco}, which EtymDB traces through Arabic \emph{shāh} \textquotedblleft{}king in chess\textquotedblright{} to Persian \emph{shāh} \textquotedblleft{}king\textquotedblright{}; the \emph{matto} of \emph{scacco matto} is often linked to a phrase saying the king is beaten.
\end{cousinweb}

\begin{grammarlens}[title={What you've built}]
One more word for telling, explaining and describing at B1.
\end{grammarlens}

\subsection*{Your turn}
\begin{itemize}
\raggedright
  \item \emph{Say it:} \emph{gli scacchi}
  \item \emph{Say it:} \emph{gli scacchi}, once more
  \item \emph{Say it:} a sentence of your own with \emph{gli scacchi}
  \item \emph{Recall:} say the Italian for trip, then the Italian for to explore, then say \emph{gli scacchi} again
\end{itemize}

\begin{tcolorbox}[breakable,colback=teal!4,colframe=teal!35!black,title={Before you move on}]
What does gli scacchi mean? (\textquotedblleft{}Chess\textquotedblright{}.) Say it once more.
\end{tcolorbox}
\section[la dama]{la dama — draughts, checkers; lady}

\label{lesson:IT-C503-dama}



\begin{quote}
Before the new one: say the Italian for to explore, then the Italian for chess.
\end{quote}

\subsection*{You'll want to know: la dama}
\textbf{la dama} — \textquotedblleft{}draughts, checkers; lady\textquotedblright{}.

\begin{quote}
Giochiamo a dama dopo cena?
\end{quote}

\begin{quote}
\emph{Shall we play draughts after dinner?}
\end{quote}

\emph{La dama} is also a lady, or a partner at a dance.

\begin{cousinweb}
\emph{dama} was borrowed from French \emph{dame} \textquotedblleft{}lady\textquotedblright{} (EtymDB), which is inherited from Latin \emph{domina} \textquotedblleft{}mistress of the house\textquotedblright{}.
\end{cousinweb}

\begin{grammarlens}[title={What you've built}]
One more word for telling, explaining and describing at B1.
\end{grammarlens}

\subsection*{Your turn}
\begin{itemize}
\raggedright
  \item \emph{Say it:} \emph{la dama}
  \item \emph{Say it:} \emph{la dama}, once more
  \item \emph{Say it:} a sentence of your own with \emph{la dama}
  \item \emph{Recall:} say the Italian for to explore, then the Italian for chess, then say \emph{la dama} again
\end{itemize}

\begin{tcolorbox}[breakable,colback=teal!4,colframe=teal!35!black,title={Before you move on}]
What does la dama mean? (\textquotedblleft{}Draughts, checkers; lady\textquotedblright{}.) Say it once more.
\end{tcolorbox}
\section[il dado]{il dado — dice; stock cube}

\label{lesson:IT-C503-dado}



\begin{quote}
Before the new one: say the Italian for chess, then the Italian for draughts.
\end{quote}

\subsection*{You'll want to know: il dado}
\textbf{il dado} — \textquotedblleft{}dice; stock cube\textquotedblright{}.

\begin{quote}
Tira il dado: tocca a te.
\end{quote}

\begin{quote}
\emph{Throw the dice: it is your turn.}
\end{quote}

The plural is \emph{i dadi}. In the kitchen, \emph{il dado} is a stock cube.

\begin{cousinweb}
The origin of \emph{dado} is not settled; it is often traced to Latin \emph{datum} \textquotedblleft{}something given\textquotedblright{}, from \emph{dare} \textquotedblleft{}to give\textquotedblright{}.
\end{cousinweb}

\begin{grammarlens}[title={What you've built}]
One more word for telling, explaining and describing at B1.
\end{grammarlens}

\subsection*{Your turn}
\begin{itemize}
\raggedright
  \item \emph{Say it:} \emph{il dado}
  \item \emph{Say it:} \emph{il dado}, once more
  \item \emph{Say it:} a sentence of your own with \emph{il dado}
  \item \emph{Recall:} say the Italian for chess, then the Italian for draughts, then say \emph{il dado} again
\end{itemize}

\begin{tcolorbox}[breakable,colback=teal!4,colframe=teal!35!black,title={Before you move on}]
What does il dado mean? (\textquotedblleft{}Dice; stock cube\textquotedblright{}.) Say it once more.
\end{tcolorbox}
\section[il gioco da tavolo]{il gioco da tavolo — board game}

\label{lesson:IT-C503-gioco-da-tavolo}



\begin{quote}
Before the new one: say the Italian for draughts, then the Italian for dice.
\end{quote}

\subsection*{You'll want to know: il gioco da tavolo}
\textbf{il gioco da tavolo} — \textquotedblleft{}board game\textquotedblright{}.

\begin{quote}
Il sabato sera facciamo un gioco da tavolo con gli amici.
\end{quote}

\begin{quote}
\emph{On Saturday evenings we play a board game with friends.}
\end{quote}

The plural is \emph{i giochi da tavolo}.

\begin{cousinweb}
\emph{gioco da tavolo} is built from \emph{gioco} \textquotedblleft{}game\textquotedblright{}, inherited from Latin \emph{iocus} \textquotedblleft{}a joke, a jest\textquotedblright{} (EtymDB), and \emph{tavolo} \textquotedblleft{}table\textquotedblright{}, from Latin \emph{tabula} \textquotedblleft{}a board\textquotedblright{} (Lewis \& Short).
\end{cousinweb}

\begin{grammarlens}[title={What you've built}]
One more word for telling, explaining and describing at B1.
\end{grammarlens}

\subsection*{Your turn}
\begin{itemize}
\raggedright
  \item \emph{Say it:} \emph{il gioco da tavolo}
  \item \emph{Say it:} \emph{il gioco da tavolo}, once more
  \item \emph{Say it:} a sentence of your own with \emph{il gioco da tavolo}
  \item \emph{Recall:} say the Italian for draughts, then the Italian for dice, then say \emph{il gioco da tavolo} again
\end{itemize}

\begin{tcolorbox}[breakable,colback=teal!4,colframe=teal!35!black,title={Before you move on}]
What does il gioco da tavolo mean? (\textquotedblleft{}Board game\textquotedblright{}.) Say it once more.
\end{tcolorbox}
\section[barare]{barare — to cheat (at games)}

\label{lesson:IT-C503-barare}



\begin{quote}
Before the new one: say the Italian for dice, then the Italian for board game.
\end{quote}

\subsection*{You'll want to know: barare}
\textbf{barare} — \textquotedblleft{}to cheat (at games)\textquotedblright{}.

\begin{quote}
Non gioco più con te: bari sempre!
\end{quote}

\begin{quote}
\emph{I am not playing with you any more: you always cheat!}
\end{quote}

Built on \emph{il baro} \textquotedblleft{}a card cheat\textquotedblright{}. In exams, \emph{copiare}.

\begin{cousinweb}
\emph{barare} is built on \emph{baro} \textquotedblleft{}a cheat at cards\textquotedblright{} (EtymDB); the origin of \emph{baro} is not settled.
\end{cousinweb}

\begin{grammarlens}[title={What you've built}]
One more word for telling, explaining and describing at B1.
\end{grammarlens}

\subsection*{Your turn}
\begin{itemize}
\raggedright
  \item \emph{Say it:} \emph{barare}
  \item \emph{Say it:} \emph{barare}, once more
  \item \emph{Say it:} a sentence of your own with \emph{barare}
  \item \emph{Recall:} say the Italian for dice, then the Italian for board game, then say \emph{barare} again
\end{itemize}

\begin{tcolorbox}[breakable,colback=teal!4,colframe=teal!35!black,title={Before you move on}]
What does barare mean? (\textquotedblleft{}To cheat (at games)\textquotedblright{}.) Say it once more.
\end{tcolorbox}
